\section*{Appendix B. Recorded infrastructure continuation}
The original outer scheduler and its child controllers received SIGTERM after three campaigns had finished; the signal's sender was not identified. Six running campaigns were affected: our harness at seeds 42/44, DrugEvolve at 42/43/44, and direct generation at 44. Completed campaigns were not restarted. Before final-test access, the operator snapshotted the original dispatches, budgets, and transport receipts and specified a common continuation rule. Originally queued campaigns ran unchanged in an independent background supervisor.

A complete recorded response can be replayed only when the original prompt and response identities match; replay does not issue another model call. A dead process with neither a complete answer nor a completed-turn event permits one byte-identical prompt retry in the same unfinished slot. The interrupted invocation and the retry both consume the original call ceiling; unknown token usage remains explicitly missing. The rule does not depend on score, add candidate slots, refund failures, or substitute seeds. Completed source/seed/contract-verified fits and native sampling decisions are reused rather than rerun or selected retrospectively.

External continuation entrypoints restore the recorded stage, native incumbent, and remaining axes or slots. They have their own hashes and receipts; the nineteen original frozen runtime files remain unchanged. For future DrugEvolve sampling after process loss, the Python RNG is reset once to that campaign's original seed; the already sampled current parent is preserved, and no alternative RNG state is searched. This is a disclosed continuation rather than a bitwise counterfactual replay. It is operator-managed infrastructure recovery, not a claim that the original research controllers autonomously recover from crashes. All affected campaigns remain in the reported result set and costs.
